\chapter{The unified theorem: Conjecture 5.10 at $k = 3, 4, 5$}
\label{ch:unified}

Aboulker, Aubian, Charbit and Lopes (arXiv:2310.04265) conjecture
(their Conjecture 5.10) that for every $k \ge 3$ there are infinitely
many $k$-$\bar\omega$-critical tournaments, and ask (their Question
5.9) whether a tournament with $\bar\omega(T) \ge k$ must contain a
\emph{bounded-size} subtournament with $\bar\omega \ge k$. The library
settles every humanly known case of the conjecture and refutes the
question at each such $k$.

\begin{trustbox}
\noindent\textbf{Trust.} Every theorem in this chapter is verified by
the Rocq kernel and reports
\rocqinline@Closed under the global context@
(\rocqinline@Print Assumptions@) --- no axioms, no admitted steps. The
witnesses are fully explicit: $\mathrm{AC}_n$,
$\mathrm{AC}_n[C_3]$, $\mathrm{AC}_n[\mathrm{AC}_n]$
(\cref{def:AC,def:T4family,def:T5family}). The small instances were
independently cross-checked by exact and SAT-based computation
(\texttt{scripts/}).
\end{trustbox}

\resulthead{The headline: every $k \in \{3,4,5\}$}{conjecture_5_10_at_345}

\begin{lstlisting}[language=Rocq]
Theorem conjecture_5_10_at_345 (k : nat) : (3 <= k <= 5)%N ->
  forall N : nat, exists T : tournament, kcritical k T /\ (N < #|T|)%N.
\end{lstlisting}
\rocqsource{theories/applications/unified.v}{66}


\begin{decoded}
For every $k$ with $3 \le k \le 5$ and every bound $N$ there exists a
$k$-$\bar\omega$-critical tournament on more than $N$ vertices: the
$k$-critical tournaments form an infinite family. (``$T$ is
$k$-critical'' unfolds per \cref{def:kcritical}: $\bar\omega(T) = k$
and $\bar\omega(T - v) = k - 1$ for every vertex $v$.)
\end{decoded}

\input{../web/panels/conjecture_5_10_at_345}

\resulthead{$k = 3$: infinitely many, and Question 5.9 fails}{conjecture_5_10_at_k3}

\begin{lstlisting}[language=Rocq]
Theorem conjecture_5_10_at_k3 (N : nat) :
  exists T : tournament, kcritical 3 T /\ (N < #|T|)%N.
\end{lstlisting}
\rocqsource{theories/applications/unified.v}{33}

\begin{lstlisting}[language=Rocq]
Theorem question_5_9_fails_at_k3 (L : nat) :
  exists T : tournament,
    [/\ kcritical 3 T, (L < #|T|)%N
      & forall S : {set T}, (3 <= `$\bar\omega$`(sub_tournament S))%N -> S = [set: T]].
\end{lstlisting}
\rocqsource{theories/applications/unified.v}{47}

\label{res:question_5_9_fails_at_k3}

\begin{decoded}
First statement: for every $N$ there is a 3-critical tournament on
more than $N$ vertices (the witness is $\mathrm{AC}_n$ itself,
$n = 2\max(2,N){+}3$). Second: moreover the witness can be chosen so
that \emph{its only subtournament with $\bar\omega \ge 3$ is the whole
tournament} --- formally, any vertex subset $S$ whose induced
subtournament has $\bar\omega \ge 3$ must be all of $V$
(\rocqinline@[set: T]@ is the full set). Hence no function
$\ell(3)$ as in Question 5.9 can exist: witnesses have unbounded size
but no proper witnessing subtournament at all.
\end{decoded}

\input{../web/panels/question_5_9_fails_at_k3}

\resulthead{$k = 4$}{conjecture_5_10_at_k4}

\begin{lstlisting}[language=Rocq]
Theorem conjecture_5_10_at_k4 (N : nat) :
  exists T : tournament, kcritical 4 T /\ (N < #|T|)%N.
\end{lstlisting}
\rocqsource{theories/applications/k4/k4_main.v}{74}

\begin{lstlisting}[language=Rocq]
Theorem question_5_9_fails_at_k4 (L : nat) :
  exists T : tournament,
    [/\ kcritical 4 T, (L < #|T|)%N
      & forall S : {set T}, (4 <= `$\bar\omega$`(sub_tournament S))%N -> S = [set: T]].
\end{lstlisting}
\rocqsource{theories/applications/k4/k4_main.v}{89}

\label{res:question_5_9_fails_at_k4}

\begin{decoded}
As for $k = 3$, with witness $\mathrm{AC}_n[C_3]$ on $3n$ vertices
(\cref{def:T4family}; the underlying values $\bar\omega = 4$,
$\bar\omega(T - v) = 3$ are \cref{res:omegabar_T4,res:omegabar_T4_del}).
\end{decoded}

\input{../web/panels/question_5_9_fails_at_k4}

\resulthead{$k = 5$}{conjecture_5_10_at_k5}

\begin{lstlisting}[language=Rocq]
Theorem conjecture_5_10_at_k5 (N : nat) :
  exists T : tournament, kcritical 5 T /\ (N < #|T|)%N.
\end{lstlisting}
\rocqsource{theories/applications/k5/main.v}{108}

\begin{lstlisting}[language=Rocq]
Theorem question_5_9_fails_at_k5 (L : nat) :
  exists T : tournament,
    [/\ kcritical 5 T, (L < #|T|)%N
      & forall S : {set T}, (5 <= `$\bar\omega$`(sub_tournament S))%N -> S = [set: T]].
\end{lstlisting}
\rocqsource{theories/applications/k5/main.v}{127}

\label{res:question_5_9_fails_at_k5}

\begin{decoded}
As above, with witness $\mathrm{AC}_n[\mathrm{AC}_n]$ on $n^2$
vertices (\cref{def:T5family}).
\end{decoded}

\input{../web/panels/question_5_9_fails_at_k5}
